\chapter{总结与展望}
\section{总结}
随着 AI 技术的不断发展和广泛应用，人们对于 AI 也越来越依赖。然而对于 AI 系统，人们却很少理解其背后的复杂供应链是否真正安全。本文围绕 AI 系统供应链的安全问题，将 AI 系统拆解成软件应用层、模型框架层和硬件加速层，并逐层分析它们的安全性问题。然而分析每一层的安全问题都会面临独有的挑战，这包括软件应用层的供应链依赖图规模庞大且复杂，模型框架层的代码实现复杂，包括跨语言调用和大量抽象以及模型供应链的不透明性；最后还有硬件加速层的供应链黑盒不开源的特点。本文为每一层分别设计了创新性地安全分析框架，以解决这些特定的分析挑战。此外本文在每一层不仅发现了新型的攻击面，更是实现了特定的扫描、检测、逆向工具来实现对该层的防护。总的来说，本文的主要工作可以分为以下三个方面：

\textbf{（1）提出了 ModuleGuard 框架实现对软件应用层供应链中模块冲突问题的安全分析与检测。} 在软件应用层，本文关注了一个新型的模块冲突问题，该问题比传统的依赖冲突问题更为细粒度，并且隐藏于包管理机制和模块解析机制中，斌没有被之前的工作注意到。本文实现了 ModuleGuard 模块冲突分析框架，并对生态中的开源软件包和项目进行大规模冲突检测。该框架不仅证实了模块冲突问题的普遍性，更是向开发者和项目维护者报告了大量此类问题，得到了他们的肯定和修复。此外本文还发现了一类新型的模块替换攻击，该攻击利用依赖包安装时同名模块会覆盖的特性对 AI 系统造成严重威胁。

\textbf{（2）提出了针对模型框架层 TensorAbuse 攻击威胁的自动化分析和恶意模型扫描工具。} 在模型框架层，本文发现了一个基于 TensorAbuse 攻击的新型的恶意模型构造方式，不同于传统的依赖漏洞或者篡改参数等对框架层的攻击，该攻击采用了合法功能 API 的滥用的方式。本文不仅构建了自动化分析框架系统性地分析了 TensorFlow 框架中可被滥用的 API 列表及其能力，并基于框架结果，构造了能够实现 TensorAbuse 攻击的攻击原语，并将攻击原语嵌入现实的 Yamnet 模型中实现真实恶意模型，上传至开源模型库。本文揭示了现有开源模型库并无法扫描此类恶意模型的现象，并实现了一套恶意模型扫描工具来防御此类攻击。

\textbf{（3）提出了针对硬件加速层的 NVIDIA GPU ASLR 逆向分析工具。}
在硬件加速层，本文专门逆向了 AI 市场占有率极高的 NVIDIA GPU，分析其上实现的安全机制是否真的安全。具体来说由于缺乏公开文档和源码，本文设计和实现了一套半自动化的 NVIDIA GPU ASLR 逆向分析工具，不仅首次分析了 GPU 上的地址空间布局和 ASLR 的实现机制，还发现了 GPU 上的 ASLR 实现非常薄弱，其仅仅实现了部分随机，甚至随机化的部分也是共享同一个随机偏移。此外本文还首次揭示了 GPU 与 CPU 地址空间之间存在的强关联性，并展示如何利用这些特定构造跨设备的攻击，从而打破 CPU 侧的 ASLR 机制保护。

\section{展望}
尽管本文已经以分层分析的方式将 AI 系统拆解成三个层级，并对每一层及的供应链单独分析，发现了多个不同层级的风险，也实现了多种创新的技术和框架来实现对各个层级的扫描、检测和保护，然而现有工作仍然存在不少局限性。

首先是针对软件应用层的模块冲突检测问题，目前的分析工作仍然是根据静态的信息抽取和模拟的方式实现的，对于软件包的动态安装行为仍然兼容性较弱。而对于这种动态性强的包，后续工作可以利用 LLM 的代码理解能力来实现更细粒度的代码逻辑分析，而不是采用静态分析的方式，从而提高分析模块信息和依赖信息的准确性。此外，在和厂商交互的过程中，他们很大一部分都不愿意修复相关问题，认为这是依赖包的维护者的问题，这会增加他们的工作量，所以修复也存在明显的滞后性，因此后续工作可以实现自动修复的工具，而不是仅仅是检测工具，从而可以减少开发者的工作量。

其次是针对模型框架层的 TensorAbuse 攻击模型的构造问题，目前的攻击构造主要是针对 TensorFlow 框架展开，虽然方法和框架具备一定的通用性，但是还没做拓展到其他 PyTorch 等框架中，因此后续工作可以拓展到更多的 AI 框架上，并分析更多的序列化格式可能存在的安全问题。此外，目前针对的恶意模型的工作都是以攻击和检测为主，后续工作也可以实现对于这类恶意模型的防护工具，例如沙箱隔离等，来防止模型造成的系统级攻击。此外，模型框架层中还有 AI 框架本身的漏洞还没有被挖掘，尤其是大模型导致衍生了大量 AI 基础设施的漏洞，如何自动化挖掘和分析这些漏洞也是未来的研究方向。

最后是针对模型框架层的 NVIDIA GPU ASLR 分析工作，目前工作只针对于 NVIDIA GPU 的特定架构和特定版本而言，在更多的硬件平台上的泛化能力仍然不足，启发式的逆向方法并没法迁移至 AMD 等计算平台。此外 GPU 上除了 ASLR 机制外，还实现了大量的安全保护机制，例如堆栈金丝雀，这些安全机制是否存在缺陷仍然不得而知，因此后续工作可以对 NVIDIA GPU 或者更多计算平台的 GPU 上的各种安全机制进行逆向分析。还有就是针对 GPU 侧的防御措施仍然较少，目前利用防护薄弱的 GPU 来实现针对 CPU 侧的攻击的防护措施也是未来的热点方向。